\documentclass[11pt,twoside]{amsart}
\usepackage{amssymb, amsmath, enumerate, palatino, hyperref}
\usepackage[normalem]{ulem}
\usepackage{fullpage}
\usepackage[T1]{fontenc}
\renewcommand{\labelitemi}{\guillemotright}
\usepackage{mathrsfs}

\usepackage{tikz}
\tikzstyle{ball} = [circle,shading=ball, ball color=black,
    minimum size=1mm,inner sep=1.3pt]

\theoremstyle{plain}
\newtheorem{prop}{Proposition}%[section]
\newtheorem{lemma}[prop]{Lemma}
\newtheorem{thm}[prop]{Theorem}
\newtheorem{obs}[prop]{Observation}
\newtheorem{app}[prop]{Application}
\newtheorem*{MainThm}{Main Theorem}
\newtheorem{cor}[prop]{Corollary}
\newtheorem{conj}[prop]{Conjecture}
\theoremstyle{remark}
\newtheorem{rmk}[prop]{Remark}
\newtheorem{prob}{Problem}
\newtheorem{bonus}[prop]{Bonus Problem}
\newtheorem{exc}{Exercise}
\theoremstyle{definition}
\newtheorem{ex}[prop]{Example}
\theoremstyle{definition}
\newtheorem{defn}[prop]{Definition}

\newcommand{\RR}{\mathbb{R}}
\newcommand{\ZZ}{\mathbb{Z}}
\newcommand{\CC}{\mathbb{C}}
\newcommand{\NN}{\mathbb{N}}
\newcommand{\QQ}{\mathbb{Q}}
\newcommand{\PP}{\mathbb{P}}

\newcommand{\cC}{\mathscr{C}}
\newcommand{\Ob}{\operatorname{Ob}}
\newcommand{\Mor}{\operatorname{Mor}}

\newcommand{\Set}{\operatorname{Set}}
\newcommand{\FinSet}{\operatorname{FinSet}}
\newcommand{\Vect}{\operatorname{Vect}}
\newcommand{\FinVect}{\operatorname{FinVect}}
\newcommand{\Mat}{\operatorname{Mat}}
\newcommand{\Gp}{\operatorname{Gp}}
\newcommand{\FinGp}{\operatorname{FinGp}}
\newcommand{\AbGp}{\operatorname{AbGp}}
\newcommand{\Ring}{\operatorname{Ring}}
\newcommand{\CommRing}{\operatorname{CommRing}}
\newcommand{\Field}{\operatorname{Field}}
\newcommand{\Top}{\operatorname{Top}}
\newcommand{\Aut}{\operatorname{Aut}}
\newcommand{\id}{\operatorname{id}}
\DeclareRobustCommand{\stir}{\genfrac\{\}{0pt}{}}


\newcommand{\defeq}{:=}

\title{Math 113: Discrete Structures\\ Homework 24}
%\author{} %Remove the leading % and put your name here if using this .tex file as a template for your homework.

\begin{document}
\maketitle

\noindent
\textbf{Due:} Friday, April 3 at 10pm.
\bigskip

\begin{prob}
  A subset~$X$ is chosen uniformly at random from the set~$[n]$. (The word
    ``uniform'' here means that each subset is equally likely.) 
  \begin{enumerate}[(a)]
    \item What is the probability that~$X$ has an even number of
      elements?
    \item Suppose~$n\geq2$.  What is the probability that~$X$ contains~$1$ and~$n$?
    \item Suppose~$n\geq10$.  What is the probability that the smallest number in~$X$ is~$10$?
  \end{enumerate}
\end{prob}

\begin{prob}
  Suppose a bag contains balls numbered~$1,2,\dots,10$.  Choose two balls from
  the bag.
  \begin{enumerate}[(a)]
    \item What is the probability the first ball is~$5$ and the second
      is~$3$ if the ball numbered~$5$ is not put back into the bag before
      drawing the second ball?
    \item What is the probability the first ball is~$5$ and the second
      is~$3$ if the ball numbered~$5$ is put back into the bag before
      drawing the second ball?
  \end{enumerate}
\end{prob}

\begin{prob} Recall that a standard poker deck has four suits and 13 cards within each suit. The deck contains 4 aces (one of every suit).
   \begin{enumerate}[(a)]
     \item What is the probability that a five-card poker hand contains exactly
       one ace?
     \item What is the probability that a five-card poker hand contains at least
       one ace?
   \end{enumerate}
\end{prob}
\end{document}
